% AY2015 材料力学 〔I〕（転記）
% 出典: 04_材料力学/original/04_材料力学/2015/2015za.pdf
% 要約: 片持ちはり（固定端B, 自由端A, 直径dの円形断面）。自由端Aに, はり軸線に
%       垂直・水平方向に固定された剛体棒（長さL/2）の先端へ鉛直下向きに外力P。
%       縦弾性係数E, 横弾性係数G。支持反力/モーメント/トルク, SFD/BMD,
%       固定端B断面の点Cの曲げ応力・せん断応力, 点Cのモール円・主応力・最大せん断応力。
% rem: 原本は3次元の俯瞰図。ここでは軸構成（x:軸方向, y:上, z:水平横）を保った
%      簡略立体図で再現する。

\section*{〔I〕}

図1に示す片持ちはりの自由端 A に, 図に示すようにはりの軸線に垂直方向に固定された
剛体棒（水平方向）の先端に, 鉛直方向下向きに外力 $P$ が加わる場合について
次の問い (1)\,$\sim$\,(4) に答えよ. ただし, はりの縦弾性係数は $E$, 横弾性係数は $G$,
断面形状は直径 $d$ の円形とする.

\begin{center}
\begin{tikzpicture}[>=Latex]
  % はり（円柱）本体: x方向（A:左 自由端 → B:右 固定端）
  \draw (1.0,0.4) -- (7.0,0.4);
  \draw (1.0,-0.4) -- (7.0,-0.4);
  \draw (1.0,0) ellipse (0.22 and 0.4);   % 自由端Aの断面
  % 固定壁（右, 斜めスラブ）
  \fill[pattern=north east lines] (7.0,-1.3) -- (7.0,1.3) -- (8.0,1.7) -- (8.0,-0.9) -- cycle;
  \draw (7.0,-1.3) -- (7.0,1.3) -- (8.0,1.7) -- (8.0,-0.9) -- cycle;
  % 軸 (A点)
  \draw[->] (1.0,0) -- (2.2,0) node[above]{$x$};
  \draw[->] (1.0,0) -- (1.0,1.3) node[left]{$y$};
  \draw[->] (1.0,0) -- (0.35,0.55) node[above left]{$z$};
  \node[below left] at (1.0,-0.2) {A};
  \node at (6.7,0.75) {B};
  % 剛体棒（Aから下方へ, 長さL/2）と荷重P
  \draw[line width=1.4pt] (1.0,-0.1) -- (1.55,-1.5);
  \node at (0.9,-1.0) {$L/2$};
  \draw[->,very thick] (1.55,-1.5) -- (1.55,-2.6) node[below]{$P$};
  % 寸法 L（はり長さ）
  \draw[<->] (1.0,1.55) -- (7.0,1.55) node[midway,fill=white]{$L$};
  % 断面（直径d）インセット（左）
  \begin{scope}[shift={(-1.3,-0.2)}]
    \draw (0,0) circle (0.5);
    \draw[->] (0,-0.5) -- (0,0.9) node[left]{$y$};
    \draw[->] (0,0) -- (-1.0,0) node[above]{$z$};
    \draw[<->] (-0.5,-0.75) -- (0.5,-0.75) node[midway,fill=white,font=\scriptsize]{$d$};
  \end{scope}
  % B断面インセット（右下, 点C）
  \begin{scope}[shift={(7.3,-2.4)}]
    \draw (0,0) circle (0.5);
    \draw[dash dot] (-0.8,0) -- (0.8,0);
    \draw[dash dot] (0,-0.8) -- (0,0.8);
    \fill (0,-0.5) circle (1.6pt);
    \node[right] at (-0.9,0.5) {C};
    \draw[->] (-0.75,0.45) -- (-0.05,-0.45);
    \node at (0,-1.1) {B 断面};
  \end{scope}
\end{tikzpicture}

\vspace{1mm}
図1
\end{center}

\begin{enumerate}[label=(\arabic*)]
  \item 支持力, 支持モーメント, 支持トルクを図に示したうえで求めよ.

  \item せん断力図（SFD）および曲げモーメント図（BMD）を求めよ.

  \item 固定端 B の円形断面の点 C に生じる曲げ応力およびせん断応力を求めよ.

  \item 固定端 B の円形断面の点 C におけるモールの応力円を描き,
        主応力および最大せん断応力を求めよ.
\end{enumerate}
